\documentclass[12pt, a4paper]{article}

\usepackage[utf8]{inputenc}
\usepackage[spanish]{babel}
\usepackage{amsthm, amsmath, amssymb}
\usepackage{hyperref}
\usepackage{csquotes}
\usepackage{biblatex}
\addbibresource{bibliography.bib}
\addbibresource{roadmap.bib}
\usepackage{geometry}
\usepackage{booktabs}
\usepackage{enumitem}
\usepackage{xcolor}
\usepackage{array}
\usepackage{longtable}

\geometry{margin=2.5cm}
\setlength{\parskip}{0.8em}
\setlength{\parindent}{0em}

\newcommand{\gate}[1]{\text{#1}}
\newcommand{\stateket}[3]{\ket{#1}^{#2}_{#3}}
\newcommand{\stategate}[3]{\gate{#1}^{#2}_{#3}}
\newcommand{\ket}[1]{|#1\rangle}
\newcommand{\bra}[1]{\langle #1|}
\newcommand{\braket}[2]{\langle #1|#2\rangle}
\newcommand{\pro}[1]{\textbf{#1}}
\newcommand{\con}[1]{\textit{#1}}
\newcommand{\gap}[1]{\textcolor{red}{\textbf{[Laguna: #1]}}}

\title{%
  \textbf{Análisis Crítico del Framework ALRCC}\\
  \large Utilidad, Ámbito, Alcance y Necesidad del Lenguaje Algebraico\\
  para Representación de Circuitos Cuánticos
}
\author{Análisis interno de investigación}
\date{Mayo 2026}

\begin{document}

\maketitle
\tableofcontents
\newpage

% ====================================================================
\section{Introducción y contexto}
% ====================================================================

El proyecto ALRCC (\textit{Algebraic Language for Representation of Quantum Circuits})
propone una notación algebraica para circuitos cuánticos que extiende la notación
bra-ket estándar mediante índices y conjuntos de índices.
Parte de una publicación previa del grupo~\cite{reina2024ix} y la amplía con el
tratamiento basado en conjuntos para manejar un número arbitrario de índices.

Los elementos notacionales núcleo son:

\begin{itemize}
  \item \textbf{Estado indexado:} $\stateket{a}{k}{\gamma}$ denota el ket
        $\ket{a}$ con el qubit $k$ reemplazado por el valor $\gamma$.
  \item \textbf{Puerta localizada:} $\gate{A}_k$ aplica la puerta $\gate{A}$
        exclusivamente al qubit $k$.
  \item \textbf{Puerta sobre conjunto:} $\gate{A}_\Gamma$ aplica $\gate{A}$
        a cada qubit del conjunto $\Gamma \subseteq \mathbb{N}_n$.
  \item \textbf{Puerta controlada general:} $\stategate{A}{\Gamma}{\Lambda}$
        con conjunto de controles $\Gamma$ y conjunto de targets $\Lambda$,
        $\Gamma \cap \Lambda = \emptyset$.
\end{itemize}

Este documento recoge un análisis crítico del estado actual del trabajo.

% ====================================================================
\section{Utilidad}
% ====================================================================

\subsection{Aportaciones reales}

\begin{enumerate}
  \item \textbf{Compacidad paramétrica.}
        La escritura $\gate{X}_\Gamma$ para aplicar una puerta a un conjunto
        arbitrario de qubits es genuinamente legible y escala con $n$ sin
        aumentar la longitud de la expresión. Esto es ventajoso en algoritmos
        donde el número de qubits es un parámetro (QFT, sumadores, redes de
        acarreo, fan-out de control).

  \item \textbf{Identidad algebraica de conjuntos.}
        La proposición
        \[
          \gate{A}_\Gamma \gate{A}_\Lambda
          = \gate{A}_{\Gamma \cup \Lambda}\, \gate{A}_{\Gamma \cap \Lambda}
        \]
        es una pieza realmente \emph{reescribible}: dada una expresión con
        dos aplicaciones de la misma puerta sobre conjuntos solapados, la
        identidad permite simplificarla en un solo paso. Éste es el tipo de
        resultado que justifica un lenguaje algebraico.

  \item \textbf{Patrón mecánico de prueba para puertas controladas.}
        La descomposición de $\stategate{A}{\Gamma}{\Lambda}$ en función de
        $a_{h_1 0}/a_{h_1 1}$ proporciona una estrategia de prueba
        reproducible y reutilizable para demostrar propiedades de circuitos
        con múltiples controles.
\end{enumerate}

\subsection{Limitaciones de utilidad}

\begin{enumerate}
  \item \textbf{Ausencia de conjunto de reglas de reescritura cerrado.}
        Los resultados existentes son lemas dispersos, no un
        \emph{cálculo} en sentido formal. Un usuario del lenguaje no dispone
        de un conjunto definido de movimientos legales, y por tanto no puede
        realizar simplificaciones sistemáticas.

  \item \textbf{Ausencia de forma normal.}
        Sin una forma normal, no es posible verificar la equivalencia de dos
        circuitos por reducción sintáctica, ni comparar el coste antes y
        después de una optimización. El archivo \texttt{sum\_generic.tex}
        (344 líneas de expansión exhaustiva) es el síntoma más claro: el
        lenguaje no está siendo usado para manipulación algebraica; se está
        usando para cómputo manual indexado.

  \item \textbf{Sin implementación operativa.}
        Sin un parser o una librería (Python, Mathematica, Lean) que ejecute
        reescrituras, el lenguaje no alcanza la práctica del campo.
\end{enumerate}

% ====================================================================
\section{Ámbito}
% ====================================================================

\subsection{Qué cubre actualmente}

\begin{longtable}{p{4cm} p{4cm} p{4.5cm}}
  \toprule
  \textbf{Cubierto}                      & \textbf{Parcialmente}                               & \textbf{Ausente}               \\
  \midrule
  Puertas 1-qubit con índice             & Controles y targets múltiples                       & Mediciones
  (aparece $_{{\Gamma}}\rfloor$ en \texttt{sum.tex} sin definición)                                                             \\
  \addlinespace
  Conmutación en qubits disjuntos        & Hadamard, Pauli, $\gate{R}_X/\gate{R}_Y/\gate{R}_Z$ &
  Operadores no unitarios, canales, ruido                                                                                       \\
  \addlinespace
  $\gate{A}_\Gamma$ con conjunto general & Ejemplos Deutsch y sumador                          & Estados
  entrelazados generales                                                                                                        \\
  \addlinespace
  Involución de puertas auto-inversas    & ---                                                 & Algoritmos con entrelazamiento
  esencial (Grover, QFT, Shor)                                                                                                  \\
  \bottomrule
\end{longtable}

\subsection{El problema del entrelazamiento}

Éste es el punto de ámbito más serio.
La descomposición
\[
  \ket{a} = a_{k0}\, \stateket{a}{k}{0} + a_{k1}\, \stateket{a}{k}{1}
\]
y la notación $a_{k\gamma} = \braket{a_k}{\gamma}$ como amplitud del $k$-ésimo
qubit están \textbf{bien definidas únicamente en estados producto}
$\ket{a} = \ket{a_1} \otimes \cdots \otimes \ket{a_n}$.

Un estado de Bell $\frac{1}{\sqrt{2}}(\ket{00}+\ket{11})$ no admite esas
amplitudes monoqubit sin ambigüedad: $a_{1,0}$ no está definido independientemente.
Mientras esto no se trate de forma explícita, la semántica del lenguaje es
válida para estados separables más superposiciones controladas, pero
\textbf{no para estados cuánticos generales}.

Para extender el lenguaje al caso general se necesitaría reemplazar
$a_{k\gamma}$ por el proyector parcial
$(\bra{\gamma}_k \otimes I_{\text{resto}})\ket{a}$, con la semántica
correspondiente.

% ====================================================================
\section{Alcance}
% ====================================================================

\subsection{Poder expresivo}

El lenguaje es, en la práctica, sub-Clifford+T: soporta puertas de Pauli,
Hadamard, rotaciones $\gate{R}_A(\theta)$ y controladas, pero la aritmética de
fase global y los estados no-stabilizer no son objetos del cálculo sino valores
de los coeficientes $a_{k\gamma}$.

\subsection{Poder de prueba}

Las proposiciones demostradas hasta ahora son, en su mayoría:
\begin{itemize}
  \item \con{Triviales por estructura tensorial:} conmutación de puertas en
        qubits distintos (inmediata porque $\gate{A} \otimes I$ e $I \otimes \gate{B}$
        siempre conmutan).
  \item \con{Elementales:} involución ($\gate{A}^2 = \gate{I}$ si $\gate{A}$
        es auto-inversa), propiedad de eigenvalores $-1$.
\end{itemize}

\textbf{No existe todavía un resultado no trivial demostrado en el lenguaje
  que sea difícil con la notación matricial estándar.}
Ese resultado sería el primer test de valor real del formalismo.

\subsection{Sin soundness ni completitud}

No hay enunciado ni demostración de que:
\begin{itemize}
  \item Dos expresiones iguales en el lenguaje corresponden a la misma
        matriz unitaria (\textbf{soundness}).
  \item Dos matrices iguales pueden reducirse a la misma expresión en el
        lenguaje (\textbf{completitud}).
\end{itemize}
Sin estos teoremas, el lenguaje no tiene estatuto de \emph{cálculo}.

% ====================================================================
\section{Necesidad}
% ====================================================================

\subsection{El problema del estado del arte ausente}

La bibliografía actual (46 entradas) no contiene ninguna referencia a:
ZX-calculus, Coecke, Selinger, Backens, Vilmart, Penrose graphical calculus,
Quipper, Silq, OpenQASM 3, PBS-calculus.

Solo se recogen~\cite{Hutsell} y~\cite{Staton} en la línea de representaciones
algebraicas de circuitos. Esto es una laguna grave: cualquier revisor de una
revista de computación cuántica devolverá el manuscrito sin leer la sección
técnica si la introducción no contextualiza frente a estas alternativas.

\subsection{Las alternativas existentes}

\begin{description}[leftmargin=2cm, style=nextline]
  \item[ZX-calculus~\cite{CoeckeDuncan2008, Backens2014, JeandelvPerdrix2018}]
        Lenguaje gráfico con semántica categórica. Demostrado completo para Clifford
        (Backens 2014), para Clifford+T (Jeandel-Perdrix-Vilmart 2018) y para
        circuitos cuánticos generales (Vilmart 2019). Tiene herramientas maduras
        (PyZX~\cite{KissingerWetering2019}, Quantomatic) y resultados de reducción
        de T-count publicados. \textbf{Es la propuesta de referencia para
          manipulación algebraica de circuitos.}

  \item[Tensor network notation / Penrose]
        Notación algebraica con índices contraídos; ya cubre lo que
        $\gate{A}_\Gamma$ expresa y conecta con simulación mediante MPS/TT.

  \item[OpenQASM 3~\cite{CrossOQASM2022}]
        Lenguaje textual ejecutable con tipos cuánticos y soporte paramétrico.

  \item[Quipper~\cite{GreenQuipper2013}]
        DSL en Haskell para circuitos cuánticos de alto nivel con generación
        paramétrica y optimización.

  \item[Silq~\cite{BichselSilq2020}]
        Primer lenguaje cuántico de alto nivel con descomputación segura e
        inferencia de tipos.

  \item[Lambda-cálculo cuántico~\cite{SelingerValiron2006}]
        Fundamento teórico para lenguajes cuánticos tipados con
        semántica denotacional completa.
\end{description}

\subsection{La pregunta que el manuscrito debe responder}

\begin{quote}
  \textit{¿Qué problema concreto resuelve ALRCC que ZX-calculus, tensor
    networks u OpenQASM no resuelven mejor?}
\end{quote}

Tres hipótesis de posicionamiento, ordenadas por defensibilidad:

\begin{description}[leftmargin=2cm, style=nextline]
  \item[(H1) Circuitos clásico-reversibles en criptografía cuántica.]
        ALRCC sería especialmente adecuada para circuitos Toffoli-based
        (sumadores, multiplicadores, funciones hash reversibles), donde la
        estructura de control-conjunto $\stategate{A}{\Gamma}{\Lambda}$ es
        natural. El ejemplo del sumador binario (\texttt{sum.tex}) apunta
        directamente en esta dirección. Esta hipótesis es defendible y conecta
        con la línea del grupo~\cite{reina2024ix}.
        \textbf{Si se acepta esta hipótesis, el alcance honesto del manuscrito
          es ``circuitos cuánticos clásico-reversibles'', no ``circuitos cuánticos
          en general''.}

  \item[(H2) Notación pedagógica / algebraica accesible.]
        ALRCC es más próxima al bra-ket clásico y por tanto más accesible para
        estudiantes de criptografía cuántica que quieren introducir optimización
        simbólica sin aprender la teoría de categorías que subyace a ZX.

  \item[(H3) Sintaxis lineal en \LaTeX / texto plano.]
        ZX requiere diagramas; ALRCC se escribe en línea. Sin embargo,
        OpenQASM ya ocupa este nicho en la práctica.
\end{description}

% ====================================================================
\section{Problemas técnicos concretos}
% ====================================================================

\begin{enumerate}
  \item \textbf{Polisemia del tercer argumento de $\stateket{a}{k}{\cdot}$.}
        El subíndice tercer-argumento puede ser un bit ($0$, $1$), una puerta
        ($\gate{A}$) o una mezcla ($1\gate{A}$ en puertas controladas). Esta
        ambigüedad viola la condición básica de que un lenguaje formal tenga
        tipos sintácticos distinguibles.

  \item \textbf{Errata en la base ortonormal.}
        En \texttt{notation.tex} y \texttt{qubit.tex} se escribe
        $\ket{+} = \frac{\ket{0}+\ket{1}}{2}$.
        La normalización correcta es $\ket{+} = \frac{\ket{0}+\ket{1}}{\sqrt{2}}$.

  \item \textbf{Medición sin definición.}
        El operador $_{{\Gamma}}\rfloor\ket{\cdot}$ aparece en \texttt{sum.tex}
        sin definición, sin propiedades y sin prueba de coherencia con el
        postulado de proyección. Las mediciones son imprescindibles para
        cualquier algoritmo cuántico.

  \item \textbf{Secciones incompletas.}
        \texttt{universal\_gates.tex} está prácticamente vacío y
        \texttt{deutsch.tex} trunca el cálculo en el momento crítico
        (ecuación~(2) inacabada). Si el lenguaje no puede completar el
        algoritmo de Deutsch-Jozsa, el ejemplo paradigmático del paper está
        roto.

  \item \textbf{La expansión de \texttt{sum\_generic.tex} contradice el objetivo.}
        El fichero de 344 líneas muestra una expansión exhaustiva que no aplica
        ninguna de las propiedades algebraicas demostradas. Si el propósito del
        lenguaje es la \emph{manipulación algebraica}, la demostración del
        sumador debería ser corta; no larga.

  \item \textbf{Proposiciones triviales elevadas a resultados.}
        La conmutación de puertas en qubits distintos ($\gate{A}_k \gate{A}_h
          = \gate{A}_h \gate{A}_k$ para $k \neq h$) es consecuencia inmediata
        de la factorización tensorial, no un resultado del lenguaje.
        Presentarla como proposición con prueba sugiere escasez de resultados
        no triviales.

  \item \textbf{Mezcla de idiomas.}
        El manuscrito alterna entre inglés y español (párrafos en
        \texttt{notation.tex}, \texttt{cgate.tex} y \texttt{deutsch.tex}),
        lo que impide su envío a cualquier revista o conferencia internacional.

  \item \textbf{Contenido duplicado.}
        \texttt{qubit.tex} y \texttt{notation.tex} contienen el mismo bloque
        sobre la notación bra-ket.
\end{enumerate}

% ====================================================================
\section{Recomendaciones prioritarias}
% ====================================================================

Las recomendaciones se ordenan por impacto en la publicabilidad del trabajo.

\begin{enumerate}[label=\textbf{R\arabic*.}]

  \item \textbf{Declarar el nicho de forma honesta y citarlo desde el
          inicio.}
        Optar por la hipótesis H1: ``lenguaje algebraico para optimización
        simbólica de circuitos clásico-reversibles en contexto criptográfico,
        como complemento textual a ZX-calculus''. Esto hace el alcance
        defendible y la contribución original.

  \item \textbf{Añadir una sección de estado del arte que cite y compare
          con ZX-calculus.}
        Como mínimo: Coecke-Duncan 2008, Backens 2014,
        Jeandel-Perdrix-Vilmart 2018, Kissinger-van de Wetering 2019 (PyZX),
        OpenQASM 3, Quipper.
        Sin esto el paper no supera la revisión editorial.

  \item \textbf{Formalizar el cálculo.}
        Definir BNF de la sintaxis, semántica denotacional (functor a
        matrices unitarias), y un conjunto de reglas de reescritura. Enunciar
        y demostrar al menos \emph{soundness}. El objetivo a medio plazo es
        un fragmento de completitud para circuitos Toffoli-based.

  \item \textbf{Demostrar un resultado palanca.}
        Encontrar un teorema de optimización (reducción de puertas Toffoli,
        conteo de CNOTs, simplificación de carry-lookahead) que el lenguaje
        demuestre en $k$ pasos frente a $k'$ pasos en la notación estándar,
        con $k \ll k'$. Ese es el argumento central del paper.

  \item \textbf{Extender la semántica al entrelazamiento.}
        Redefinir $a_{k\gamma}$ como proyector parcial para estados generales,
        o delimitar explícitamente el dominio a estados separables.

  \item \textbf{Añadir definición y propiedades de la medición.}
        Formalizar $_{{\Gamma}}\rfloor\ket{\cdot}$, su semántica probabilística
        y la reducción de estado post-medición.

  \item \textbf{Implementar un prototipo.}
        Un módulo Python/SymPy o un notebook Mathematica que realice
        reescrituras del lenguaje y exporte a Qiskit. Es lo que hace adoptable
        una propuesta de notación.

  \item \textbf{Correcciones de manuscrito.}
        \begin{itemize}
          \item Corregir la normalización de $\ket{+}$ (factor $1/\sqrt{2}$).
          \item Traducir todos los párrafos en español.
          \item Completar o eliminar \texttt{universal\_gates.tex} y
                \texttt{deutsch.tex}.
          \item Eliminar la duplicación entre \texttt{qubit.tex} y
                \texttt{notation.tex}.
          \item Resolver la polisemia del tercer argumento de
                $\stateket{a}{k}{\cdot}$ mediante una gramática formal.
        \end{itemize}
\end{enumerate}

% ====================================================================
\section{Línea de trabajo: papers a leer y orden recomendado}
% ====================================================================

La línea de trabajo se estructura en cuatro fases. El criterio de orden es
``necesidad mínima antes de continuar escribiendo'': cada fase desbloquea
los argumentos que la siguiente requiere.

% ----- FASE 1 -------------------------------------------------------
\subsection*{Fase 1 — Fundamentos y estado del arte imprescindible (2--3 semanas)}

\textit{Objetivo: entender qué existe ya y dónde encaja ALRCC.}

\begin{enumerate}
  \item[\checkmark] \textbf{Coecke \& Duncan (2008) --- Interacting quantum observables.}
        Artículo fundacional del ZX-calculus. Establece los spiders, las
        reglas de reescritura y la semántica. Leer completo.
        \textit{Por qué primero:} es la propuesta más directamente competidora;
        todo el discurso de posicionamiento de ALRCC necesita conocerla.

  \item \textbf{Backens (2014) --- The ZX-calculus is complete for stabilizer quantum mechanics.}
        Primer resultado de completitud de ZX para el fragmento Clifford.
        \textit{Por qué:} da el estándar que ALRCC debería aspirar a igualar
        en su nicho.

  \item \textbf{Kissinger \& van de Wetering (2019) --- PyZX: Large scale automated diagrammatic reasoning.}
        Herramienta práctica que implementa reducción de T-count usando
        ZX-calculus.
        \textit{Por qué:} caso de uso concreto de optimización con un lenguaje
        algebraico; muestra el \emph{gap} que ALRCC necesita cerrar.

  \item \textbf{Hutsell \& Greenwood (2009) --- Efficient algebraic representation of quantum circuits.}
        Ya está en la bibliografía~\cite{Hutsell}. Leer en profundidad como
        antecedente directo de ALRCC.

  \item[\checkmark] \textbf{Nielsen \& Chuang (2010), Capítulos 4--5.}
        Revisar la notación estándar de circuitos, puertas universales y la
        gate por la que ALRCC extiende el formalismo. Ya en la
        bibliografía~\cite{NielsenChuang}.
\end{enumerate}

% ----- FASE 2 -------------------------------------------------------
\subsection*{Fase 2 --- Fundamentos de cálculos de reescritura (2 semanas)}

\textit{Objetivo: adquirir las herramientas formales para definir el cálculo de ALRCC.}

\begin{enumerate}[resume]
  \item \textbf{Baader \& Nipkow (1998) --- Term Rewriting Systems.}
        Capítulos 1--3. Definición de sistemas de reescritura de términos (TRS),
        terminación, confluencia, formas normales.
        \textit{Por qué:} para formalizar las reglas de reescritura de ALRCC
        y probar que el cálculo termina y es confluente.

  \item \textbf{Abramsky \& Coecke (2004) --- A categorical semantics of quantum protocols.}
        Introducción a la mecánica cuántica como categoría compacta.
        \textit{Por qué:} ZX-calculus vive aquí; conocer el marco categórico
        mínimo permite comparar ALRCC a ese nivel y decidir si vale la pena
        una formulación categórica.

  \item \textbf{Selinger (2004) --- Towards a quantum programming language.}
        Primer diseño sistemático de un lenguaje de programación cuántico con
        semántica denotacional.
        \textit{Por qué:} modelo a seguir para dar semántica formal a ALRCC.
\end{enumerate}

% ----- FASE 3 -------------------------------------------------------
\subsection*{Fase 3 --- Optimización y aplicaciones a circuitos reversibles (2--3 semanas)}

\textit{Objetivo: identificar el resultado palanca (R4) e implementar el prototipo.}

\begin{enumerate}[resume]
  \item \textbf{Jeandel, Perdrix \& Vilmart (2018) --- A complete axiomatisation of the ZX-calculus for Clifford+T.}
        Completitud de ZX para Clifford+T. Leer las reglas del sistema: son el
        modelo a seguir para definir el sistema de reglas de ALRCC para
        circuitos reversibles.

  \item \textbf{Amy, Maslov, Mosca \& Roetteler (2013) --- A meet-in-the-middle algorithm for fast synthesis of depth-optimal quantum circuits.}
        Algoritmo de síntesis y optimización de circuitos Clifford+T.
        \textit{Por qué:} métrica concreta de optimización (profundidad,
        conteo de T) que ALRCC podría mejorar para el fragmento Toffoli.

  \item \textbf{Maslov (2016) --- Advantages of using relative-phase Toffoli gates with an application to multiple control Toffoli optimization.}
        Resultado directo sobre optimización de puertas Toffoli multi-control.
        \textit{Por qué:} ALRCC con $\stategate{X}{\Gamma}{\Lambda}$ debería
        poder expresar y mejorar estos resultados algebraicamente.

  \item \textbf{Cross et al. (2022) --- OpenQASM 3: A broader and deeper quantum assembly language.}
        Especificación completa de OpenQASM 3.
        \textit{Por qué:} para posicionar ALRCC frente a un lenguaje textual
        ejecutable existente y detectar la diferenciación real.
\end{enumerate}

% ----- FASE 4 -------------------------------------------------------
\subsection*{Fase 4 --- Entrelazamiento, mediciones y extensión del alcance (3--4 semanas)}

\textit{Objetivo: resolver las lagunas de ámbito y publicar una versión ampliada.}

\begin{enumerate}[resume]
  \item \textbf{Selinger \& Valiron (2006) --- A lambda calculus for quantum computation with classical control.}
        Lambda-cálculo cuántico tipado con semántica denotacional completa.
        \textit{Por qué:} modelo para extender ALRCC a mediciones y control clásico.

  \item \textbf{Bichsel et al. (2020) --- Silq: A high-level quantum language with safe uncomputation and its type system.}
        Lenguaje cuántico con descomputación segura automática.
        \textit{Por qué:} la descomputación es esencial en circuitos
        reversibles con ancillas; ALRCC necesita modelarla.

  \item \textbf{Perdrix \& Wang (2016) --- Supplementarity is necessary for quantum diagram reasoning.}
        Demuestra que algunas propiedades de ZX son irreducibles.
        \textit{Por qué:} indica qué resultados de completitud son
        necesariamente difíciles y dónde ALRCC podría tener ventaja al
        restringir el dominio.

  \item \textbf{Green et al. (2013) --- Quipper: A scalable quantum programming language.}
        Diseño de Quipper: generación paramétrica de circuitos cuánticos
        a gran escala.
        \textit{Por qué:} caso de estudio de lenguaje cuántico adoptado
        en la práctica; modelo para el prototipo de ALRCC.
\end{enumerate}

% ----- RESUMEN -------------------------------------------------------
\subsection*{Mapa de lectura resumido}

\begin{center}
  \begin{tabular}{clll}
    \toprule
    \textbf{Semana} & \textbf{Paper clave}  & \textbf{Objetivo}         & \textbf{Desbloquea}     \\
    \midrule
    1--2            & Coecke-Duncan 2008    & Conocer el competidor     & R2 (estado del arte)    \\
    1--2            & Backens 2014          & Estándar de completitud   & R3 (formalizar cálculo) \\
    2               & Kissinger-W. 2019     & Caso de uso concreto      & R4 (resultado palanca)  \\
    2               & Hutsell 2009          & Antecedente directo       & R2                      \\
    3               & Baader-Nipkow cap.1-3 & Herramientas TRS          & R3                      \\
    3               & Selinger 2004         & Semántica formal          & R3                      \\
    4--5            & Jeandel-P-V 2018      & Reglas Clifford+T         & R3, R4                  \\
    5               & Maslov 2016           & Métrica Toffoli           & R4                      \\
    5--6            & Amy et al. 2013       & Optimización de circuitos & R4                      \\
    6               & OpenQASM 3 2022       & Diferenciación textual    & R1                      \\
    7--8            & Selinger-Valiron 2006 & Mediciones tipadas        & R6                      \\
    7--8            & Bichsel et al. 2020   & Descomputación            & R6                      \\
    9               & Green et al. 2013     & Modelo prototipo          & R7                      \\
    \bottomrule
  \end{tabular}
\end{center}

% ====================================================================
\section{Veredicto}
% ====================================================================

\begin{description}[leftmargin=2cm, style=nextline]
  \item[Utilidad]
        Real para circuitos clásico-reversibles parametrizados por conjuntos;
        no demostrada fuera de ese nicho.

  \item[Ámbito]
        Sub-Clifford+T en la práctica; estados generales entrelazados y
        mediciones quedan fuera del lenguaje en su estado actual.

  \item[Alcance]
        Notación con propiedades locales interesantes pero sin cálculo cerrado,
        sin resultado palanca y sin herramienta que justifique la adopción
        frente a ZX-calculus o OpenQASM.

  \item[Necesidad]
        \textbf{Condicionada.}
        Si el posicionamiento es ``lenguaje algebraico textual para optimizar
        circuitos reversibles en contexto criptográfico, complementario a
        ZX-calculus'', la necesidad es \textbf{defendible} y el trabajo tiene
        lugar en la literatura. Si la pretensión es competir con ZX como cálculo
        general de circuitos cuánticos, la necesidad \textbf{no está demostrada}
        y el camino exige completitud, mediciones y herramientas.
\end{description}

La recomendación más urgente es ejecutar R1 (acotar el nicho) y R2 (cubrir el
estado del arte) antes de invertir en más demostraciones internas, porque sin
esos dos movimientos el manuscrito corre el riesgo de ser rechazado sin
revisión técnica.

% ====================================================================
\printbibliography[title={Referencias --- Estado del arte a incorporar}]
% ====================================================================

\end{document}
